\documentclass[11pt,a4paper]{article}
\usepackage{brainteaser}

\begin{document}
\brainteasertitle{3}{Arrêt optimal avec un dé}

\begin{statementbox}
On part d'un total $S=0$. À chaque étape, on peut s'arrêter et recevoir $S$
euros, ou relancer un dé équilibré.

Si le dé donne $1$, la partie s'arrête immédiatement et le gain est nul. Si le
dé donne $2,3,4,5$ ou $6$, le résultat est ajouté à $S$, puis on choisit à
nouveau entre s'arrêter et continuer.

Déterminer la stratégie optimale, le seuil d'arrêt et le gain moyen optimal en
partant de $S=0$.
\end{statementbox}

\begin{hintsbox}
\begin{enumerate}
  \item Définir $V(S)$, la valeur optimale lorsque le total courant vaut $S$.
  \item Écrire l'équation de Bellman
  \[
  V(S)=\max\left(S,\frac16\sum_{k=2}^{6}V(S+k)\right).
  \]
  \item Si l'on relance une fois puis s'arrête après un résultat non fatal,
  quelle espérance obtient-on?
  \item Au-dessus d'un certain seuil, montrer que le score forme un
  supermartingale tant que l'on continue.
  \item Une fois le seuil trouvé, calculer $V(19),V(18),\ldots,V(0)$ à rebours.
\end{enumerate}
\end{hintsbox}

\correction

Si l'on relance une fois puis s'arrête après un résultat non fatal, l'espérance
est
\[
\frac16\sum_{k=2}^{6}(S+k)
=\frac{5S+20}{6}
=S+\frac{20-S}{6}.
\]
Pour $S<20$, cette quantité est strictement supérieure à $S$: continuer au
moins une fois est préférable à l'arrêt immédiat.

Pour $S\geq20$, le score devient $0$ après un $1$, et $S+k$ après un autre
résultat. Son espérance au coup suivant est inférieure ou égale à $S$. Après
un lancer non fatal, le score reste au moins $20$. Le score arrêté est donc un
supermartingale, et aucune stratégie de continuation ne peut avoir une
espérance supérieure au gain certain $S$.

\begin{keyideabox}
La stratégie optimale est
\[
\boxed{\text{continuer si }S<20,\qquad\text{s'arrêter si }S\geq20.}
\]
À $S=20$, les deux décisions ont la même espérance; s'arrêter évite le risque.
\end{keyideabox}

Pour calculer la valeur initiale, on pose $V(S)=S$ pour $S\geq20$, puis, pour
$0\leq S<20$,
\[
V(S)=\frac16\sum_{k=2}^{6}V(S+k).
\]
Le calcul à rebours donne notamment
\[
V(19)=\frac{115}{6},\qquad
V(18)=\frac{55}{3},\qquad
V(0)=\frac{492303203}{60466176}.
\]

\begin{resultbox}
Le seuil d'arrêt est $20$ et le gain moyen optimal en partant de zéro vaut
\[
\boxed{V(0)\simeq 8{,}142\text{ euros}.}
\]
\end{resultbox}

\end{document}
